\begin{abstract}
  % Equal-contribution note. The title block is set in a box (see main.tex) and
  % footnotes cannot escape from a box, so the marker is written by hand there
  % and the note itself is emitted here, from ordinary page body.
  \blfootnote{$^{*}$These authors contributed equally to this work.}%
  PDF translation is a joint machine-translation and reconstruction problem:
  pages store positioned glyphs rather than semantic paragraphs, and translated
  text rarely occupies the same visual space as the source. This paper studies
  where errors arise in a layout-preserving PDF translation pipeline and which of
  them its design can control. We hypothesise that document structure analysis is
  a binding stage because its errors are unrecoverable downstream, and that
  layout fidelity is better computed during reconstruction than requested from a
  language model as a length constraint. Our three-stage system parses every page
  visually with layout detection, OCR and table-cell detection, translates
  addressable elements with an identifier-preserving LLM contract, and
  reconstructs the page by estimating font size and appearance from the source
  geometry and pixels. Stage-level evaluation supports this design. On 1{,}651
  OmniDocBench pages, the parser reaches $F_1@.5 = 0.940$ under a
  granularity-invariant matcher and $0.889$ label accuracy; symbol recognition is
  strong (formula CDM $F_1 = 0.863$), so the remaining loss is one of
  \emph{detection}---formulas on three-column pages and roughly a fifth of
  tables---rather than reading. The same shape appears at cell level on
  PubTables-1M, used because OmniDocBench has no per-cell boxes: $F_1 = 0.902$
  overall, but recall of $0.504$ on cells spanning several rows or columns.
  On all 55 \texttt{en}$\rightarrow$\texttt{xx}
  WMT24++ pairs, five LLMs span $0.788$--$0.841$ COMET-DA, with smaller
  between-model than between-language variation. Finally, a $\pm15\%$
  character-length constraint is honoured only 69--73\% of the time and is
  uncorrelated with adequacy ($r = 0.066$). These results are not a substitute
  for an end-to-end PDF benchmark, but they identify the stages and document
  types where such a benchmark should focus.
\end{abstract}
